% overview-chapter.tex — 总览章骨架（v2 教学化版式，无习题）
% 用法：整板验证通过后，基于全工程产物填实本骨架；与模块章共用 main-demo.tex 导言区。
% 内容四件套：指标追溯表 / 电源树与信号流 / 板级布局策略与模块间接口 / 整体验证汇总。

\chapter{系统总览：从需求到制造文件}

\begin{introduction}
  \item 对照需求追溯表，指出每项硬约束由哪个模块、哪条验证证据支撑（分析）；
  \item 画出（或复述）整机电源树与信号流（理解）；
  \item 陈述板级布局策略：功能分区、分地策略与单点连接（理解）；
  \item 汇总各模块验证结果，指出哪些结论为 VERIFIED、哪些仍为 COMPUTED（记住）。
\end{introduction}

\section{需求与指标追溯}

⟨引言：整机要解决什么问题、交付形态（层数/尺寸/接口）。⟩

⟨指标追溯表：每行一项硬约束，列 = 需求编号 / 需求内容 / 承接模块 / 设计措施 / 验证方法 / 证据位置（reports/…）/ 状态徽标。这一张表是"许愿→兑现"的总账，逐行必须能回链 MISSION.md 的 H1…Hn 与 reports/ 下的记录。⟩

\begin{center}
\begin{tabular}{@{}lllll@{}}
  \toprule
  需求 & 内容 & 承接模块 & 验证方法与证据 & 状态 \\
  \midrule
  H1 & ⟨…⟩ & ⟨module-id⟩ & ⟨ERC/断言/仿真/实测 + reports/ 路径⟩ & ⟨状态⟩ \\
  \bottomrule
\end{tabular}
\end{center}

\section{电源树与信号流}

⟨整机电源树：输入→各级变换→负载，标注电压/预计最大电流/纹波敏感负载；用 circuitikz 或 TikZ 重绘（教学重绘版），复杂处拆子图。⟩

⟨信号流：传感→调理→ADC→MCU→执行/显示，标注模拟段与数字段的分界——它决定分地与分区。⟩

\begin{proposition}{分地与单点连接}{ground-strategy}
⟨本板分地策略：功率地/模拟地划分依据，0Ω 桥接位置与理由；两面铺地与散热孔回链。规则来源：PCB 培训 + 本工程决策日志条目。⟩
\end{proposition}

\section{板级布局策略与模块间接口}

⟨功能分区图（TikZ 框图或布局截图标注）：每区放什么模块、为什么相邻/远离；敏感网络与电感/天线的避让关系。⟩

⟨模块间接口表：接口名 / 连接双方 / 网络 / 载流或带宽要求 / 布线要求（线宽/差分/开尔文）。IR 中各模块 boundary_ports 的汇总视图。⟩

\section{整体验证计划与结果汇总}

⟨验证计划表：层级（模块/整板）/ 验证手段（ERC、网表断言、DRC 铺铜前后、独立 Gerber 复核、SPICE 或计算复核、冷重启复测）/ 判据 / 结果 / 证据路径。全部结论标注状态徽标；实测未做的如实标未建立。⟩

⟨遗留与二期：打样后实测项清单（指向各模块章的实测占位表）。⟩

\begin{conclusion}
⟨3--5 句：整机方案一句话；追溯表证明"每条需求都有出处与证据"；布局与分地策略的贯穿思想；当前证据链中 VERIFIED 与 COMPUTED 的边界——明确告诉读者哪些数字要等打样实测才能定稿。⟩
\end{conclusion}
